% =====================================================================
%  TP1 — Presentación de defensa (10 min + 8 min de preguntas)
%  72.75 Aprendizaje Automático — ITBA — 2026 Q2
%
%  Compilar:  cd informe && pdflatex presentacion.tex && pdflatex presentacion.tex
%
%  Sobre las animaciones: las que importan son los OVERLAYS (revelado
%  progresivo), que funcionan en cualquier visor de PDF. Las transiciones
%  de página (\transdissolve y similares) sólo las respetan Adobe Reader e
%  Impress!ve; en Preview, Skim o el visor de Chrome no hacen nada. Por eso
%  el peso de la animación está puesto en los overlays y las transiciones
%  van como agregado.
% =====================================================================
\documentclass[aspectratio=169,11pt]{beamer}

\usepackage[T1]{fontenc}
\usepackage[utf8]{inputenc}
\usepackage[spanish,es-noquoting]{babel}
\usepackage{lmodern}
\usepackage{amsmath,amssymb}
\usepackage{graphicx}
\usepackage{booktabs}
\usepackage{tikz}
\usetikzlibrary{arrows.meta,positioning,shapes.geometric,fit,backgrounds,decorations.pathreplacing}

\graphicspath{{../figuras/}{../figuras/presentacion/}}

% ---------------------------------------------------------------------
% Tema: mínimo y construido a mano.
% Los temas que trae Beamer (Madrid, Warsaw) gastan una franja entera de
% pantalla en barras de navegación que nadie usa en una defensa de 10
% minutos. Acá la pantalla es para los datos.
% ---------------------------------------------------------------------
\usetheme{default}
\usecolortheme{default}
\setbeamertemplate{navigation symbols}{}
\setbeamertemplate{itemize items}[circle]
\setbeamertemplate{blocks}[rounded][shadow=false]

% Paleta: la misma que las figuras, validada para separar bajo daltonismo.
\definecolor{azul}{HTML}{2A78D6}
\definecolor{naranja}{HTML}{EB6834}
\definecolor{azuloscuro}{HTML}{104281}
\definecolor{tinta}{HTML}{0B0B0B}
\definecolor{tintasec}{HTML}{52514E}
\definecolor{apagado}{HTML}{898781}
\definecolor{rejilla}{HTML}{E1E0D9}
\definecolor{superficie}{HTML}{FCFCFB}

\setbeamercolor{normal text}{fg=tinta,bg=superficie}
\setbeamercolor{structure}{fg=azul}
\setbeamercolor{frametitle}{fg=tinta,bg=superficie}
\setbeamercolor{title}{fg=tinta}
\setbeamercolor{alerted text}{fg=naranja}
\setbeamercolor{block title}{fg=superficie,bg=azul}
\setbeamercolor{block body}{fg=tinta,bg=rejilla!35}
\setbeamerfont{frametitle}{size=\large,series=\bfseries}
\setbeamerfont{title}{size=\LARGE,series=\bfseries}

% Título de slide con una regla fina debajo: separa sin gritar.
\setbeamertemplate{frametitle}{%
  \vspace{0.6em}\hspace{0.4em}\insertframetitle%
  \vspace{0.25em}\newline\hspace*{0.4em}%
  \textcolor{rejilla}{\rule{0.965\paperwidth}{0.8pt}}%
}

% Pie: sección a la izquierda, número de slide a la derecha. Nada más.
% El ancho es \paperwidth MENOS el margen que se agrega adentro; sumarle leftskip y
% rightskip a un box que ya mide \paperwidth lo desborda en cada slide.
\setbeamertemplate{footline}{%
  \hbox{\begin{beamercolorbox}[wd=\paperwidth,ht=2.4ex,dp=1.2ex]{}%
    \hspace*{1em}%
    {\tiny\color{apagado}\insertsectionhead}\hfill
    {\tiny\color{apagado}\insertframenumber}%
    \hspace*{1em}%
  \end{beamercolorbox}}%
}


% Slide donde la figura es la protagonista: el titulo se achica, la figura toma casi
% toda la altura util y debajo va UNA sola linea. En una defensa la gente mira el
% grafico, no lee vinetas; el texto que compite con la figura no se lee y encima la
% empequeñece.
\newcommand{\figgrande}[2][0.99]{%
  \begin{center}\includegraphics[width=#1\textwidth,height=0.80\textheight,%
    keepaspectratio]{#2}\end{center}}

% Atajos de énfasis

% Toda figura entra en una CAJA de ancho y alto, con keepaspectratio.
% Dimensionar sólo por ancho obliga a conocer el aspecto de cada imagen y falla en
% cuanto una es más alta que ancha: en 16:9 la restricción que ata es el ALTO.
% Con las dos cotas, LaTeX escala por la que ate primero y nada se desborda.
\newcommand{\figslide}[2][0.95]{%
  \includegraphics[width=#1\textwidth,height=0.66\textheight,keepaspectratio]{#2}}

% Atajos de énfasis
\newcommand{\dato}[1]{\textcolor{azul}{\textbf{#1}}}
\newcommand{\ojo}[1]{\textcolor{naranja}{\textbf{#1}}}
\newcommand{\gris}[1]{\textcolor{apagado}{#1}}

% Un número grande, para las slides donde la respuesta ES un número.
\newcommand{\heroe}[2]{%
  \begin{center}
    {\fontsize{46}{50}\selectfont\textcolor{azul}{\textbf{#1}}}\\[0.35em]
    {\small\color{tintasec}#2}
  \end{center}%
}

\title{Regresión e introducción a la evaluación de modelos}
\subtitle{Trabajo Práctico 1}
\author{72.75 --- Aprendizaje Automático}
\institute{Instituto Tecnológico de Buenos Aires}
\date{26 de agosto de 2026}

% Los numeros de TEST entran por aca. Los genera `python3 -m src.evaluar_test`; mientras
% no se haya corrido, el archivo es la plantilla y los documentos muestran [PENDIENTE].
% Asi ningun valor de test queda transcripto a mano.
% PLANTILLA. Este archivo lo SOBREESCRIBE `python3 -m src.evaluar_test`.
%
% Mientras el test no se haya evaluado, los documentos compilan igual y muestran
% "PENDIENTE" donde iria cada numero. Asi nunca hay un valor de test transcripto a
% mano en el informe: o esta el que produjo la evaluacion, o se ve que falta.
\newcommand{\pendiente}{\textcolor{naranja}{\textbf{[PENDIENTE]}}}
\newcommand{\rmsetestproduccion}{\pendiente}
\newcommand{\rmsetestganador}{\pendiente}
\newcommand{\rmsetestlineal}{\pendiente}
\newcommand{\rmsetestbaseline}{\pendiente}
\newcommand{\costosimplicidad}{\pendiente}
\newcommand{\featuresproduccion}{\pendiente}
\newcommand{\coefsproduccion}{\pendiente}
\newcommand{\featuresganador}{\pendiente}
\newcommand{\coefsganador}{\pendiente}
\newcommand{\rmsevalproduccion}{4955,28}
\newcommand{\rmsevalganador}{4919,97}
\newcommand{\rmsetestredondeado}{\pendiente}
\newcommand{\testevaluado}{no}


\begin{document}

% =====================================================================
\begin{frame}[plain]
  \titlepage
  \vspace{-1.2em}
  \begin{center}
    \small\color{tintasec}
    Dataset: \texttt{Insurance Charges} --- 1338 registros, 7 variables
  \end{center}
\end{frame}

% =====================================================================
\section{El problema}
% =====================================================================
\begin{frame}{Qué hay que predecir}
  \transdissolve
  \begin{columns}[T,onlytextwidth]
    \begin{column}{0.54\textwidth}
      \vspace{0.5em}
      Predecir el \dato{costo médico anual} de una persona a partir de:
      \vspace{0.6em}
      \begin{itemize}\setlength\itemsep{0.35em}
        \item<2-> \texttt{age}, \texttt{bmi}, \texttt{children} \gris{--- numéricas}
        \item<3-> \texttt{sex}, \texttt{smoker} \gris{--- binarias}
        \item<4-> \texttt{region} \gris{--- nominal, 4 categorías}
      \end{itemize}
      \vspace{0.9em}
      \onslide<5->{%
        \begin{block}{La pregunta real del TP}
          No es cuál modelo ajusta mejor.\\
          Es \ojo{cuál modelo elegiría} y \ojo{qué error prometería}.
        \end{block}}
    \end{column}
    \begin{column}{0.44\textwidth}
      \vspace{1.2em}
      \onslide<6->{%
      \begin{center}
      \begin{tikzpicture}[font=\small]
        \node[draw=rejilla,fill=rejilla!40,rounded corners=3pt,minimum width=3.1cm,
              minimum height=0.85cm,align=center] (x) {6 variables};
        \node[draw=azul,fill=azul!12,rounded corners=3pt,minimum width=3.1cm,
              minimum height=0.85cm,below=0.75cm of x,align=center] (m) {modelo};
        \node[draw=naranja,fill=naranja!12,rounded corners=3pt,minimum width=3.1cm,
              minimum height=0.85cm,below=0.75cm of m,align=center] (y) {\texttt{charges}};
        \draw[-{Stealth[length=2.4mm]},tintasec,thick] (x) -- (m);
        \draw[-{Stealth[length=2.4mm]},tintasec,thick] (m) -- (y);
      \end{tikzpicture}
      \end{center}}
    \end{column}
  \end{columns}
\end{frame}

% =====================================================================
\section{Train, validación y test}
% =====================================================================
\begin{frame}{Cómo se evalúa: los tres conjuntos y la validación cruzada}
  \transwipe
  \vspace{-0.1em}
  \begin{center}
  \begin{tikzpicture}[font=\scriptsize,x=1.42cm,y=0.58cm]
    % --- los tres conjuntos ---
    \fill[azul!18,rounded corners=2pt] (0,0.4) rectangle (5,1.3);
    \node[font=\small] at (2.5,0.85) {\textbf{train} — 1070 filas};
    \fill[naranja!18,rounded corners=2pt] (5.15,0.4) rectangle (6.4,1.3);
    \node[font=\small] at (5.78,0.85) {\textbf{test} — 267};
    \node[right=2mm,text=tintasec,align=left] at (6.4,0.85)
      {se toca \textbf{una vez}, al final};

    % --- los 5 folds, revelados de a uno ---
    \foreach \f in {1,...,5}{
      \onslide<\the\numexpr\f+1\relax->{
      \node[left=1mm,text=apagado] at (0,{-\f*0.95}) {fold \f};
      \foreach \b in {1,...,5}{
        \pgfmathtruncatemacro{\v}{ifthenelse(\b==\f,1,0)}
        \ifnum\v=1
          \fill[naranja,rounded corners=1.5pt]
            ({(\b-1)},{-\f*0.95-0.3}) rectangle ({\b-0.06},{-\f*0.95+0.3});
          \node[white,font=\tiny] at ({\b-0.53},{-\f*0.95}) {valida};
        \else
          \fill[azul!22,rounded corners=1.5pt]
            ({(\b-1)},{-\f*0.95-0.3}) rectangle ({\b-0.06},{-\f*0.95+0.3});
          \node[tintasec,font=\tiny] at ({\b-0.53},{-\f*0.95}) {entrena};
        \fi}}}
    \onslide<7->{
      \draw[decorate,decoration={brace,amplitude=4pt},apagado]
        (5.05,-0.85) -- (5.05,-4.9)
        node[midway,right=5pt,align=left,text=tintasec,font=\scriptsize]
        {se promedian\\los 5 errores};}
  \end{tikzpicture}
  \end{center}

  \vspace{0.2em}
  \begin{columns}[T,onlytextwidth]
    \begin{column}{0.48\textwidth}
      \onslide<8->{\footnotesize
      \textbf{Por qué no alcanza el error de train}\\[0.15em]
      El modelo ajustó sus parámetros para minimizarlo. Un valor bajo puede ser que
      aprendió \emph{o} que memorizó: no distingue.}
    \end{column}
    \begin{column}{0.48\textwidth}
      \onslide<9->{\footnotesize
      \textbf{Por qué el test se toca una sola vez}\\[0.15em]
      En cuanto se lo usa para decidir algo pasa a ser un segundo conjunto de
      validación, y su número deja de estimar el error real.}
    \end{column}
  \end{columns}
\end{frame}


% =====================================================================
\section{Limpieza de datos}
% =====================================================================
\begin{frame}{Los outliers: ¿error de carga o subpoblación real?}
  \transwipe
  \vspace{-0.2em}
  \begin{center}
    \only<1>{\figgrande[0.97]{04-outliers-charges.png}}
    \only<2->{\figgrande[0.97]{04-outliers-charges.png}}
  \end{center}
  \vspace{-0.4em}
  \onslide<2->{%
  \begin{center}
    \small
    \dato{139 outliers} por IQR (10,4\,\%) --- y el \ojo{97,8\,\%} son fumadores,
    contra 11,5\,\% en el resto.\\[0.2em]
    \color{tintasec}Un error de carga estaría repartido al azar. Esto no lo está:
    \textbf{se conservan}.
  \end{center}}
\end{frame}

% ---------------------------------------------------------------------
\begin{frame}{El criterio importa: IQR contra z-score}
  \vspace{0.4em}
  \begin{center}
  \begin{tabular}{lrr}
    \toprule
    \textbf{Variable} & \textbf{IQR detecta} & \textbf{Z-score detecta} \\
    \midrule
    \texttt{charges}  & \dato{139} & \gris{7} \\
    \texttt{children} & \gris{0}   & \dato{18} \\
    \bottomrule
  \end{tabular}
  \end{center}

  \vspace{1.0em}
  \begin{itemize}\setlength\itemsep{0.5em}
    \item<2-> En \texttt{charges} el z-score detecta 20 veces menos: la asimetría es
      1,516, y \ojo{los propios extremos inflan el desvío} que el z-score usa de
      referencia. El criterio se sabotea a sí mismo.
    \item<3-> En \texttt{children} pasa lo inverso: es un entero de 0 a 5, y
      ``tres desvíos'' no significa nada en una variable así.
  \end{itemize}
  \vspace{0.6em}
  \onslide<4->{\begin{center}\large
    Se usa \dato{IQR}, que se apoya en cuartiles y es robusto a la asimetría.
  \end{center}}
\end{frame}

% =====================================================================
\section{El pipeline}
% =====================================================================
\begin{frame}{Dentro de cada fold, en este orden}
  \vspace{0.6em}
  \begin{center}
  \begin{tikzpicture}[font=\scriptsize,node distance=3mm,
      p/.style={draw,rounded corners=3pt,minimum height=1.15cm,minimum width=2.05cm,
                align=center,thick,fill=superficie}]
    \node[p,draw=apagado] (a) {codificar\\\gris{one-hot}};
    \node[p,draw=azul,fill=azul!10,right=of a] (b) {estandarizar};
    \node[p,draw=apagado,right=of b] (c) {expandir\\\gris{polinómica}};
    \node[p,draw=azul,fill=azul!10,right=of c] (d) {estandarizar\\\gris{otra vez}};
    \node[p,draw=naranja,fill=naranja!10,right=of d] (e) {ajustar\\modelo};
    \foreach \i/\j in {a/b,b/c,c/d,d/e}
      \draw[-{Stealth[length=2mm]},tintasec] (\i) -- (\j);
  \end{tikzpicture}
  \end{center}

  \vspace{0.8em}
  \begin{itemize}\setlength\itemsep{0.45em}
    \item<2-> Los dos estandarizadores se ajustan \ojo{sólo con el sub-train del fold}.
      Ajustarlos antes de partir sería \textbf{fuga de datos}: la media de validación se
      filtraría al entrenamiento.
    \item<3-> El primer escalado evita que $\texttt{age}^3 = 262\,144$ conviva con
      $\texttt{children}^3 = 125$.
    \item<4-> El segundo hace que la penalización $L_1$ sea comparable entre features.
  \end{itemize}
  \vspace{0.5em}
  \onslide<5->{\begin{center}\small\color{tintasec}
    La regresión polinómica sigue siendo \textbf{lineal en los parámetros}:
    se transforman las columnas, no el modelo.
  \end{center}}
\end{frame}

% =====================================================================
\section{Resultados}
% =====================================================================
\begin{frame}[label=curvau]{Qué pasa al subir el grado del polinomio}
  \begin{center}
    \only<1>{\figgrande[0.90]{curva-u-1.png}}%
    \only<2>{\figgrande[0.90]{curva-u-2.png}}%
    \only<3>{\figgrande[0.90]{curva-u-3.png}}%
    \only<4->{\figgrande[0.90]{curva-u-4.png}}%
  \end{center}
  \vspace{-0.2em}
  \begin{center}\small
    \only<1>{\color{tintasec}Si mirásemos sólo esto, la conclusión sería
      ``más grado siempre es mejor''.}
    \only<2>{\color{tintasec}Mejora hasta grado 2 y después \ojo{empeora}.
      El de grado 4 es peor que el lineal simple.}
    \only<3>{\color{tintasec}La banda es $\pm 1$ desvío entre folds:
      \ojo{$\pm$244} en grado 2 contra \ojo{$\pm$1031} en grado 4.}
    \only<4->{\color{tintasec}El grado 4 no es sólo peor:
      es \ojo{inestable}, cambia según cómo caiga la partición.}
  \end{center}
\end{frame}

% ---------------------------------------------------------------------
\begin{frame}{Los dos indicadores de sobreajuste}
  \vspace{0.5em}
  \begin{columns}[T,onlytextwidth]
    \begin{column}{0.48\textwidth}
      \onslide<1->{%
      \begin{block}{1. La brecha train--validación}
        \begin{center}
          \begin{tabular}{lr}
            grado 1 & \gris{88} \\
            grado 4 & \dato{2508} \\
          \end{tabular}
        \end{center}
        {\footnotesize Un modelo que generaliza tiene los dos errores parecidos.
        Una brecha que se ensancha es la firma directa del sobreajuste.}
      \end{block}}
    \end{column}
    \begin{column}{0.48\textwidth}
      \onslide<2->{%
      \begin{block}{2. El desvío entre folds}
        \begin{center}
          \begin{tabular}{lr}
            grado 1 & \gris{$\pm$89} \\
            grado 4 & \dato{$\pm$1031} \\
          \end{tabular}
        \end{center}
        {\footnotesize Más de diez veces mayor. No dice ``peor en promedio'':
        dice \textbf{poco confiable}.}
      \end{block}}
    \end{column}
  \end{columns}
  \vspace{1.0em}
  \onslide<3->{\begin{center}
    El error de train de grado 4 es el \ojo{mejor de toda la tabla} (4058).\\[0.2em]
    \small\color{tintasec}Por eso el error de train no sirve para elegir.
  \end{center}}
\end{frame}

% ---------------------------------------------------------------------
\begin{frame}{Lasso: regularizar recupera la estabilidad}
  \vspace{-0.2em}
  \begin{center}
    \figgrande[0.99]{camino-lasso-apaisado.png}
  \end{center}
  \vspace{-0.3em}
  \onslide<1->{\begin{center}\small\color{tintasec}
    Al aflojar $\lambda$ entran más features y el error primero baja y después sube.
    \\[0.15em]
    En grado 4, la penalización $L_1$ deja vivos \dato{26 de 494} coeficientes.
  \end{center}}
\end{frame}

% =====================================================================
\section{Las tres respuestas}
% =====================================================================
\begin{frame}{1. ¿Qué modelo obtuvo menor error?}
  \vspace{1.2em}
  \heroe{Lasso grado 4}{$\lambda = 286{,}37$ \quad---\quad RMSE de validación
    $= 4920{,}0 \pm 224{,}3$}
  \vspace{0.8em}
  \onslide<1->{%
  \begin{center}
    \large\color{tintasec}
    Y sin embargo, \ojo{no es el que llevaría a producción}.
  \end{center}}
\end{frame}

% ---------------------------------------------------------------------
\begin{frame}{2. ¿Cuál implementaría en una aplicación real?}
  \begin{center}
    \only<1>{\figgrande[0.88]{una-es-1.png}}%
    \only<2->{\figgrande[0.88]{una-es-2.png}}%
  \end{center}
  \vspace{-0.3em}
  \begin{center}\small
    \only<1>{\color{tintasec}Están separadas por decenas de dólares\ldots}
    \only<2>{\color{tintasec}\ldots pero el desvío entre folds es de \ojo{224}.
      El error estándar es $224{,}3/\sqrt{5} = 100{,}3$.}
    \only<3->{\color{tintasec}\textbf{8 de 18 configuraciones son
      estadísticamente indistinguibles.}\\
      Cuál sale primera lo decide el azar de la partición, no el modelo.}
  \end{center}
\end{frame}

% ---------------------------------------------------------------------
\begin{frame}{Entre indistinguibles, el más simple}
  \vspace{0.5em}
  \begin{center}
  \begin{tabular}{lrrr}
    \toprule
     & \textbf{Features} & \textbf{Coefs. vivos} & \textbf{RMSE test} \\
    \midrule
    % Los `&` y `\\` NO pueden ir adentro de un grupo: son caracteres de alineación que
    % TeX escanea a nivel del tabular, y encerrarlos en {...} rompe el escáner (se cuelga
    % sin dar error). El revelado va celda por celda, con \visible, que además reserva el
    % espacio y evita que la tabla salte entre overlays.
    Ganador de la CV \gris{(grado 4)} & \featuresganador & \coefsganador & \rmsetestganador \\
    \visible<1->{\textbf{Producción} \gris{(grado 2)}} & \visible<1->{\dato{\featuresproduccion}}
      & \visible<1->{\dato{\coefsproduccion}} & \visible<1->{\dato{\rmsetestproduccion}} \\
    \midrule
    \visible<3->{\gris{Lineal simple (grado 1)}} & \visible<3->{\gris{8}}
      & \visible<3->{\gris{8}} & \visible<3->{\gris{\rmsetestlineal}} \\
    \visible<3->{\gris{Predecir siempre la media}} & \visible<3->{\gris{---}}
      & \visible<3->{\gris{0}} & \visible<3->{\gris{\rmsetestbaseline}} \\
    \bottomrule
  \end{tabular}
  \end{center}
  \vspace{1.0em}
  \onslide<2->{%
  \begin{center}
    La simplicidad cuesta \ojo{+92 dólares de RMSE} --- un 2\,\% ---\\
    y compra un espacio de features \dato{11 veces más chico}.
  \end{center}}
\end{frame}

% ---------------------------------------------------------------------
\begin{frame}{3. ¿Qué RMSE esperar en datos nuevos?}
  \vspace{0.6em}
  \onslide<1->{\heroe{\rmsetestproduccion}{RMSE de \textbf{test} del modelo de producción, en dólares}}
  \vspace{0.3em}
  \onslide<2->{%
  \begin{center}
    \small\color{tintasec}
    No los \textbf{4955} de validación.
  \end{center}
  \vspace{0.5em}
  \begin{block}{Por qué el número de validación está sesgado a la baja}
    \small
    Esa configuración se eligió por ser el \ojo{mínimo} de 19 estimaciones ruidosas.
    El mínimo de un conjunto de estimaciones ruidosas está sesgado hacia abajo
    \emph{aunque cada una sea insesgada}: es probable que el ganador lo sea en parte
    por una racha favorable de ruido.
  \end{block}}
  \vspace{0.3em}
  \onslide<3->{%
  \begin{center}\footnotesize\color{tintasec}
    Dos advertencias: el test tiene 267 filas (el número tiene su propia incertidumbre),
    y el RMSE promedia sobre una población heterogénea --- el error se concentra en los
    fumadores.
  \end{center}}
\end{frame}

% =====================================================================
\section{El hallazgo}
% =====================================================================
\begin{frame}{Por qué el polinomio mejora: hay dos poblaciones}
  \begin{center}
    \only<1>{\figgrande[0.86]{interaccion-1.png}}%
    \only<2>{\figgrande[0.86]{interaccion-2.png}}%
    \only<3->{\figgrande[0.86]{interaccion-3.png}}%
  \end{center}
  \vspace{-0.3em}
  \begin{center}\small
    \only<1>{\color{tintasec}Correlación de Pearson: \ojo{0,198}.
      Parece una variable poco informativa.}
    \only<2>{\color{tintasec}Eran \ojo{dos poblaciones superpuestas}.}
    \only<3->{\color{tintasec}Un factor de \ojo{17,7} entre las pendientes.
      Eso es un término $\texttt{bmi} \times \texttt{smoker}$, que un modelo
      aditivo no puede representar.}
  \end{center}
\end{frame}

% ---------------------------------------------------------------------
\begin{frame}{El Lasso lo encontró solo}
  \vspace{0.5em}
  \begin{center}
  \begin{tabular}{lr}
    \toprule
    \textbf{Feature} & \textbf{Coeficiente} \\
    \midrule
    \texttt{smoker=yes}     & 9301{,}54 \\
    \texttt{age}            & 3463{,}37 \\
    \visible<1->{\texttt{bmi*smoker=yes}} & \visible<1->{\dato{3317{,}35}} \\
    \visible<2->{\texttt{bmi}}            & \visible<2->{\gris{1595{,}81}} \\
    \bottomrule
  \end{tabular}
  \end{center}
  \vspace{0.9em}
  \onslide<2->{%
  \begin{center}
    El término de interacción quedó \ojo{por encima de \texttt{bmi} sola}.\\[0.3em]
    \small\color{tintasec}
    Sin que nadie se lo indicara, la penalización $L_1$ recuperó el mismo fenómeno
    que el análisis exploratorio había encontrado a mano.
  \end{center}}
  \vspace{0.4em}
  \onslide<3->{\begin{center}\footnotesize\color{tintasec}
    De 44 features, sobreviven 10. $L_1$ las apaga en \textbf{cero exacto};
    $L_2$ las habría encogido sin anular ninguna.
  \end{center}}
\end{frame}

% =====================================================================
\section{Cierre}
% =====================================================================
\begin{frame}{Limitaciones}
  \vspace{0.6em}
  \begin{itemize}\setlength\itemsep{0.7em}
    \item<1-> \textbf{Colinealidad en grados altos.} En grado 4, el \dato{56\,\%} de las
      494 columnas son redundantes: una dummy al cuadrado es función afín exacta de sí
      misma. Los coeficientes de un grupo colineal \ojo{no se interpretan de a uno}.
    \item<2-> \textbf{Una sola semilla.} El RMSE de test tiene varianza asociada a qué
      filas cayeron en test, que no medimos.
    \item<3-> \textbf{El dataset es simulado}, no el registro de una aseguradora real.
    \item<4-> \textbf{El RMSE penaliza al cuadrado}, así que está dominado por los casos
      caros.
  \end{itemize}
\end{frame}

% ---------------------------------------------------------------------
\begin{frame}{En una línea}
  \vspace{1.0em}
  \begin{center}
    {\large El modelo que llevaría a producción es un}\\[0.5em]
    {\Large\textcolor{azul}{\textbf{Lasso de grado 2 con 10 features vivas}}}\\[0.7em]
    {\large y esperaría un error típico de}\\[0.4em]
    {\fontsize{40}{44}\selectfont\textcolor{azul}{\textbf{\rmsetestproduccion}}}
  \end{center}
  \vspace{0.9em}
  \onslide<1->{%
  \begin{center}\small\color{tintasec}
    No es el de menor error de validación.\\
    Es el más simple entre los que son \textbf{estadísticamente indistinguibles} de él.
  \end{center}}
  \vspace{0.6em}
  \onslide<2->{\begin{center}\footnotesize\color{apagado}
    Todo implementado desde cero sobre \texttt{numpy}: split, $k$-fold, one-hot,
    estandarizado, ecuación normal, expansión polinómica y Lasso.
  \end{center}}
\end{frame}

% =====================================================================
% Slides de respaldo — para los 8 minutos de preguntas
% =====================================================================
\appendix

\begin{frame}[plain]
  \vspace{2.5cm}
  \begin{center}
    {\Large\color{tintasec}Slides de respaldo}\\[0.4em]
    {\small\color{apagado}para las preguntas}
  \end{center}
\end{frame}

\begin{frame}{Respaldo: el diagnóstico de rango}
  \vspace{0.4em}
  \begin{center}
  \begin{tabular}{crrr}
    \toprule
    \textbf{Grado} & \textbf{Columnas} & \textbf{Rango efectivo} & \textbf{Redundantes} \\
    \midrule
    1 & 8   & 8   & 0 \\
    2 & 44  & 36  & 8 \gris{(18\,\%)} \\
    3 & 164 & 100 & 64 \gris{(39\,\%)} \\
    4 & 494 & 216 & \dato{278 (56\,\%)} \\
    \bottomrule
  \end{tabular}
  \end{center}
  \vspace{0.9em}
  \begin{itemize}\setlength\itemsep{0.4em}
    \item Una dummy al cuadrado sigue teniendo dos valores:
      $\texttt{smoker}^2 = 1{,}456866 \cdot \texttt{smoker} + 1{,}0$,
      residuo $1{,}5 \times 10^{-14}$.
    \item El producto de dos dummies del mismo one-hot ya vive en el espacio generado
      por las dummies y la constante.
    \item \ojo{No} es que ese producto sea cero: lo sería en la codificación cruda 0/1,
      pero estandarizamos \emph{antes} de expandir.
  \end{itemize}
\end{frame}

\begin{frame}{Respaldo: por qué \texttt{lstsq} y no invertir $X^TX$}
  \vspace{0.6em}
  El número de condición completo de la matriz de diseño estandarizada:
  \vspace{0.5em}
  \begin{center}
  \begin{tabular}{crrrr}
    \toprule
    \textbf{Grado} & 1 & 2 & 3 & 4 \\
    \midrule
    $\kappa$ & 2{,}2 & $1{,}8\times10^{16}$ & $8{,}7\times10^{16}$
             & \dato{$3{,}1\times10^{18}$} \\
    \bottomrule
  \end{tabular}
  \end{center}
  \vspace{0.9em}
  \begin{itemize}\setlength\itemsep{0.45em}
    \item De grado 2 en adelante supera la precisión de \texttt{float64}
      ($\approx 10^{16}$): la matriz es \ojo{numéricamente singular}.
    \item \texttt{np.linalg.inv} devolvería basura. \texttt{lstsq} usa SVD y da la
      solución de norma mínima.
    \item Restringido al subespacio de rango completo el condicionamiento es benigno
      (4,0 / 18,4 / 129,4): las \textbf{predicciones} son estables. Lo que no lo es
      es el reparto de coeficientes.
  \end{itemize}
\end{frame}

\begin{frame}{Respaldo: el modelo sobre el test}
  \vspace{-0.2em}
  \begin{center}
    \IfFileExists{../figuras/05-predicho-vs-real.png}{%
      \figgrande[0.62]{05-predicho-vs-real.png}}{%
      \vspace{2cm}\large\color{tintasec}
      Esta figura se genera al correr \texttt{python -m src.evaluar\_test}.\\[0.4em]
      \normalsize Predecir sobre test \emph{es} evaluar el test, así que no se dibuja antes.}
  \end{center}
  \vspace{-0.4em}
  \begin{center}\small\color{tintasec}
    Acierta con precisión en la mayoría barata y se dispersa en los casos caros,
    que son casi todos fumadores.\\
    Es la asimetría del error que la tercera respuesta advierte.
  \end{center}
\end{frame}

\begin{frame}{Respaldo: cómo verificamos los números}
  \vspace{0.5em}
  \begin{center}
  \begin{tabular}{ll}
    \toprule
    \textbf{Chequeo} & \textbf{Resultado} \\
    \midrule
    RMSE de CV recalculado sin usar el módulo & coincide a 0,1 \\
    RMSE de test (producción y referencias)   & coincide a 0,01 \\
    Lasso contra \texttt{scipy.optimize}      & coincide a 8 decimales \\
    Usos de \texttt{X\_test} antes del punto 5 & 4 usos, 0 sospechosos \\
    Alineación de los 494 nombres de grado 4  & verificada \\
    Suites de tests                            & 3/3 en verde \\
    \bottomrule
  \end{tabular}
  \end{center}
  \vspace{0.9em}
  \begin{center}\small\color{tintasec}
    La alineación de nombres se verificó asignando un \textbf{primo distinto a cada
    columna}: el valor de cada monomio es entonces un producto de primos que factoriza
    unívocamente al monomio.
  \end{center}
\end{frame}

\end{document}
